%%================================================
%% Appendix D — Alembic migration pack
%%================================================
\chapter{ภาคผนวก ง: ลำดับการย้ายสคีมาฐานข้อมูลด้วย Alembic}
\label{app:alembic_migration_pack}

ภาคผนวกนี้สรุปภาพรวมของ migration stream ฝั่งเซิร์ฟเวอร์ที่ใช้ Alembic เป็นตัวจัดการสคีมาฐานข้อมูล เพื่อสนับสนุนความเข้าใจเชิงวิศวกรรมของบทที่~\ref{chapter3} และอธิบายพัฒนาการของโดเมนข้อมูลที่เกี่ยวข้องกับผลในบทที่~\ref{chapter4}

\section*{บทบาทของ Alembic ในแพลตฟอร์ม}

เอกสารฝั่งเซิร์ฟเวอร์ระบุชัดว่าการเปลี่ยนแปลง schema ที่กระทบระบบจริงต้องผ่าน Alembic และการทำงานของ backend อาศัยการอัปเกรด migration ให้ตรงกับสภาพโค้ดก่อนเริ่มบริการหลัก แนวทางนี้ช่วยให้ฐานข้อมูลมีประวัติการเปลี่ยนแปลงที่ตรวจสอบย้อนกลับได้ และทำให้การเพิ่มโดเมนใหม่ เช่น workflow, AI settings, chat actions, และ localization runtime เกิดขึ้นอย่างเป็นระบบ

\section*{คำสั่งหลักที่ใช้ในการดูแล migration}

คำสั่งพื้นฐานที่ใช้กับ migration stream มีดังนี้

\begin{itemize}
    \item \texttt{alembic upgrade head} สำหรับอัปเกรด schema ไปยัง revision ล่าสุด
    \item \texttt{alembic revision --autogenerate -m "..."} สำหรับสร้าง revision ใหม่จากความเปลี่ยนแปลงของ model
    \item \texttt{alembic current} สำหรับตรวจสอบ revision ปัจจุบันของฐานข้อมูล
\end{itemize}

แม้ระบบทดสอบจำนวนหนึ่งจะใช้ \texttt{create\_all()} เพื่อเร่งการตั้งค่า test database แต่เส้นทาง production และการติดตั้งจริงยังต้องอาศัย migration stream นี้เป็นหลัก

\section*{ลำดับการขยายโดเมนข้อมูลที่สำคัญ}

จากไฟล์ revision ที่อยู่ใน \texttt{server/alembic/versions/} สามารถมองเห็นแนวโน้มการขยายโดเมนหลักของระบบได้เป็นช่วงดังนี้

\begin{itemize}
    \item ช่วงตั้งต้นของระบบและ workspace/auth เช่น \texttt{7eb2ee25df34\_initial\_workspace\_schema.py}, \texttt{5f91820b6334\_add\_workspaces\_and\_dual\_mode\_.py}, และ \texttt{8957f9d4c3ba\_add\_user\_model\_for\_auth\_rbac.py}
    \item ช่วงขยายโดเมนผู้ป่วย อุปกรณ์ และ floorplan เช่น \texttt{e9f0a1b2c3d4\_phase1\_device\_management\_mvp.py}, \texttt{b7c8d9e0f1a2\_add\_floorplan\_layouts\_table.py}, และ \texttt{f1a2b3c4d5e6\_add\_device\_activity\_events.py}
    \item ช่วงขยาย workflow และงานดูแล เช่น \texttt{1aac0f420dfc\_add\_task\_management\_tables.py}, \texttt{9a6b3f4d2c10\_add\_workflow\_domain\_tables.py}, \texttt{aa1bb2cc3dd4\_add\_routine\_task\_tables.py}, และ \texttt{a1b2c3d4e5f7\_add\_unified\_task\_management.py}
    \item ช่วงขยาย AI, chat actions, MCP tokens และ runtime telemetry เช่น \texttt{e8f1a2b3c4d5\_add\_ai\_chat\_and\_workspace\_ai\_settings.py}, \texttt{p0q1r2s3t4u5\_add\_chat\_actions\_table.py}, \texttt{u5v6w7x8y9z0\_add\_mcp\_tokens\_table.py}, และ \texttt{q1r2s3t4u5v6\_add\_device\_localization\_runtime\_tables.py}
    \item ช่วงขยาย care workflow และ pipeline events เช่น \texttt{y1z2a3b4c5d6\_add\_care\_workflow\_jobs.py}, \texttt{z9\_merge\_care\_task\_wf\_job\_and\_jobs\_table.py}, และ \texttt{r2s3t4u5v6w7\_add\_pipeline\_events\_and\_behavioral\_states.py}
\end{itemize}

การมอง revision เป็นช่วงเช่นนี้ช่วยให้ผู้อ่านเห็นว่าการเติบโตของฐานข้อมูลไม่ได้เป็นเพียงการเพิ่มตารางแบบกระจัดกระจาย แต่สะท้อนวิวัฒนาการของระบบจากโครงสร้างพื้นฐานสู่แพลตฟอร์มที่รองรับงานดูแล การสื่อสาร การวิเคราะห์ และ agent runtime ในระดับที่ซับซ้อนขึ้น

\section*{สถานะของ revision head ปัจจุบัน}

จากเอกสาร runtime ของฝั่งเซิร์ฟเวอร์ มีการระบุว่าระบบในช่วงเวลานี้มี Alembic head สองตัวที่สำคัญคือ \texttt{e7f8a9b0c1d2} และ \texttt{r2s3t4u5v6w7} โดยในสภาพแวดล้อม compose สำหรับการบูตใหม่ backend ใช้คำสั่ง \texttt{alembic upgrade heads} เพื่อให้ schema ใหม่สามารถเริ่มต้นได้สำเร็จแม้ยังมีสองปลายทางของ migration stream อยู่พร้อมกัน

ประเด็นนี้มีนัยสำคัญในเชิงปฏิบัติการ เพราะสะท้อนว่าระบบวิจัยที่พัฒนาอย่างต่อเนื่องอาจมี migration branch มากกว่าหนึ่งสายในบางช่วงเวลา ผู้ดูแลระบบจึงต้องติดตามทั้ง revision graph และผลกระทบต่อการ deploy อย่างใกล้ชิด

\section*{ความสำคัญต่อการทำซ้ำผลและการบำรุงรักษา}

ผลการทดลองในบทที่~\ref{chapter4} จำนวนมากอาศัยโดเมนฐานข้อมูลที่เกิดจาก migration ภายหลัง เช่น ตาราง workflow, AI settings, chat actions, localization runtime และ pipeline events หากผู้นำระบบไปทำซ้ำใช้ schema ที่ไม่ตรงกับ revision ที่อธิบายไว้ ผลด้านฟังก์ชันและเส้นทางข้อมูลอาจเปลี่ยนไปอย่างมีนัยสำคัญ

ภาคผนวกนี้จึงทำหน้าที่เป็นแผนที่เชิงประวัติศาสตร์ของฐานข้อมูล ช่วยให้ผู้อ่านเห็นความเชื่อมโยงระหว่างการเติบโตของฟังก์ชันในแพลตฟอร์มกับการขยายสคีมาฐานข้อมูลที่รองรับฟังก์ชันเหล่านั้น
